\begin{center}\bfseries Resumen\end{center}
\phantomsection
\addcontentsline{toc}{section}{Resumen}
\begingroup
\fontsize{9.5}{10.7}\selectfont
\setlength{\parskip}{1.5pt plus .5pt minus .5pt}
% Composición local: conserva resumen y palabras clave en la apertura.
\fontsize{10.2}{11.6}\selectfont
\setlength{\parskip}{1pt plus .5pt minus .5pt}
\enlargethispage{2\baselineskip}
\noindent
Se construye una respuesta constitutiva positiva e invertible del vacío
electromagnético: dos canales orientados determinan conjuntamente las
coordenadas adimensionales de permitividad, permeabilidad, impedancia y
velocidad de propagación. Se demuestran sus relaciones exactas y la
recuperación única del par angular que las produce. La misma respuesta
permite reconstruir la forma de la elipse de acción y expresar el
funcional de Barbero--Immirzi mediante sus dos autovalores etiquetados.
El resultado relaciona así la determinación de los valores con la
información estructural conservada por su evaluación.

La construcción parte de un núcleo formal matemático denominado
Holografía Modular Triádica, en adelante HMT. \emph{All Possible Paths}
(APP) es la estructura aritmética de caminos sobre el soporte toroidal
$(\mathbb Z/9\mathbb Z)^2$; sus desplazamientos cardinales forman el
grafo de Cayley aditivo, y sus hojas conservan las evaluaciones suma y
producto con residuo y cociente. TRIT es la firma orientada
$\{-1,0,+1\}$ que distingue los regímenes locales. El
\emph{Topological Phase Kernel} (TPK) compone selección, transporte,
actualización de memoria y retorno sobre ese estado enriquecido. La
prolongación nonádica y el cierre dodecafásico generan las coordenadas
correlacionadas y los registros de los que proceden la acción, el par
angular y las incidencias utilizadas por la respuesta.

Los sectores $90$ y $120$ del transporte $T$ fijan
$\mathcal V=(I-T^{90})(I-T^{120})^{-1}$. Su función escalar
$f(s)=(1+s+s^2)/(1+s+s^2+s^3)$ es una biyección decreciente de $(0,1)$
sobre $(3/4,1)$: esta propiedad establece la inversión cúbica y su
dominio. El intercambio de hojas permuta permitividad y permeabilidad,
invierte impedancia y conserva velocidad. En la misma coordenatización dimensional
se demuestra además la identidad entre velocidad constitutiva y
velocidad cinemática del reloj.

La realización sobre el ciclo de doce fases aporta el defecto de traza
$-1/12$, una representación de $f$ como cociente de determinantes y un cuarto de giro
con cuadrado $-I$. La doble integración logarítmica de su resolvente
produce el momento de Catalán. En la sección de Pell $2-\sqrt3$ se
demuestra una identidad exacta con la constante de Catalán y
$\log(2+\sqrt3)$; los dos argumentos constitutivos mantienen su
determinación angular propia. La inversión de $f$ compone después la
respuesta del vacío con el funcional orientado de Barbero--Immirzi.
La restitución integral demuestra asimismo que la función constitutiva
determina el momento orientado completo bajo la condición
$\mathcal C_2(s)\to0$ en el origen; su límite en $s=1$ recupera
la constante de Catalán. Se conserva así una correspondencia funcional
con su orientación y sus argumentos, además de sus evaluaciones.

La polarización discreta se obtiene del balance acotado de vacancias,
y el prefactor electrónico $\sqrt3/4$ de la norma del conmutador de
las proyecciones asociadas a las fibras de suma y producto del soporte
aritmético tridimensional. La sección
positiva de carga, construida con $\alpha$, la acción y la impedancia,
determina las relaciones de von Klitzing, conductancia, flujo magnético
y Josephson, incluidas sus identidades cruzadas y su covariancia bajo
el retorno de la acción. Las realizaciones dimensionales conservan las
secciones de acción, carga y velocidad y sus cambios de unidades. La
evaluación numérica corresponde a las coordenadas internas;
su comparación física conserva explícitamente la representación y la
coordenatización dimensional utilizadas.

La separación positiva del modo común y del modo orientado especifica
qué información conserva la normalización unimodular. La resolución
de Fourier de la incidencia dodecafásica identifica cuatro modos y
sus pesos exactos, distinguiendo el defecto de rango, las multiplicidades
firmadas y el coeficiente singular del retorno.

La inversión induce asimismo una ley asociativa de composición de
respuestas, con una coordenada angular aditiva. En la realización
electrodébil de árbol, los lectores de calibre e Higgs recuperan
conjuntamente los canales y la impedancia, con dominio explícito.

\smallskip
\noindent\textbf{Palabras clave:} determinación estructural de constantes;
vacío electromagnético; respuesta constitutiva; inversión espectral;
Holografía Modular Triádica; orientación de hojas; constante de Catalán;
unidad de Pell; Barbero--Immirzi; polarización; cuantización eléctrica.
\par\endgroup
